\section{评估与证据}

\subsection{Benchmark Stack}

偏好优化的目标是让模型在多种 benchmark 上展现一致进步。常见的评估层次：

\begin{table}[H]
\centering
\begin{tabular}{p{0.2\textwidth} p{0.28\textwidth} p{0.34\textwidth}}
\toprule
层次 & 示例指标 & 说明 \\
\midrule
自动化 benchmark & HumanEval、MMLU、TruthfulQA & 量化能力与 reasoning \\
偏好一致性 & Pairwise Preference Accuracy & 奖励模型能否预测人类偏好 \\
对齐检查 & Red Team、Safety Gym & 模型在安全边界内的表现 \\
部署指标 & Rejection Rate、Hallucination Rate & 实际使用中体验质量 \\
\bottomrule
\end{tabular}
\caption{RLHF/DPO 评估栈}
\end{table}

\subsection{Evidence Matrix}

部署前需要整理 evidence matrix，搭建对齐与验证的证据链：

\begin{table}[H]
\centering
\begin{tabular}{p{0.2\textwidth} p{0.24\textwidth} p{0.24\textwidth} p{0.24\textwidth}}
\toprule
Evidence & Source & Metric & Decision Use \\
\midrule
Pairwise human ranking & crowdsourced label studio & Accuracy & Reward model calibration \\
Automated guardrails & prompt test suite & Reject rate & Safety filter tuning \\
Long-context consistency & multi-turn benchmark & Drift score & Deployment guardrail \\
Deployment logs & tooling + telemetry & Incident count & Postmortem loop \\
\bottomrule
\end{tabular}
\caption{Evidence matrix 用于跨阶段验证}
\end{table}

\subsection{Calibration 与 Robustness}

不仅需要对齐 evaluation，还要在不同语义注入、上下文长度、token budget 变动时保持稳定。常见措施有：

\begin{itemize}
    \item 对 reward model 输出做 temperature calibration 到 0-1 区间
    \item 建立 multi-context benchmark，检测 evidence drift
    \item 用 counterfactual prompt (e.g. `change numeric context`) 验证 reward 变化趋势一致
    \item 记录 rejection case IDs 以便 drill-down
\end{itemize}

\begin{warningbox}{偏好模型 drift 误判的风险}
即便 reward model 在训练集表现良好，也可能在实际 prompt 中 drift。如果没有 calibration 和 drift detection，就无法保证 alignment 在部署后仍然成立。
\end{warningbox}

\subsection{部署指标与 evidence ranking}

当模型进入灰度或全量部署后，evaluation stack 需要稳健对接监控层：
\begin{itemize}
    \item \textbf{Rejection/deferral rate}：判断 guardrail 是否过于宽松或严格。
    \item \textbf{Hallucination trend}：按 topic 分组，查看 hallucination rate 是否在特定方向聚集。
    \item \textbf{Latency drift}：策略更新后 latency 是否突破 SLO。
    \item \textbf{Evidence score}：把 benchmark、human eval、deployment logs 建立分层排序，为 decision board 提供实时依据。
\end{itemize}

\begin{knowledgebox}{Evidence ranking 的价值}
把不同阶段的 evidence 赋予不同权重（例如：部署 logs > benchmark），可以在不牺牲安全的前提下快速推进版本迭代。
\end{knowledgebox}

\begin{importantbox}{评价仪表盘需要实时可视化}
Archit 提醒：\textit{``Evidence matrix 不是写在文档里的表格，而是在 dashboard 上不断更新的热力图''}，只有实时可视化才能发现 drift 盲点。
\end{importantbox}

\subsection{本章小结}

评估与证据体系需要从 benchmark 到部署层层覆盖，Evidence Matrix 是把不同信号拼接成可执行结论的枢纽，deploy metrics 则把 evaluation stack 延伸到 production。

